% Options for packages loaded elsewhere
\PassOptionsToPackage{unicode}{hyperref}
\PassOptionsToPackage{hyphens}{url}
\PassOptionsToPackage{dvipsnames,svgnames,x11names}{xcolor}
%
\documentclass[
  11pt,
]{article}
\usepackage{amsmath,amssymb}
\usepackage{iftex}
\ifPDFTeX
  \usepackage[T1]{fontenc}
  \usepackage[utf8]{inputenc}
  \usepackage{textcomp} % provide euro and other symbols
\else % if luatex or xetex
  \usepackage{unicode-math} % this also loads fontspec
  \defaultfontfeatures{Scale=MatchLowercase}
  \defaultfontfeatures[\rmfamily]{Ligatures=TeX,Scale=1}
\fi
\usepackage{lmodern}
\ifPDFTeX\else
  % xetex/luatex font selection
\fi
% Use upquote if available, for straight quotes in verbatim environments
\IfFileExists{upquote.sty}{\usepackage{upquote}}{}
\IfFileExists{microtype.sty}{% use microtype if available
  \usepackage[]{microtype}
  \UseMicrotypeSet[protrusion]{basicmath} % disable protrusion for tt fonts
}{}
\makeatletter
\@ifundefined{KOMAClassName}{% if non-KOMA class
  \IfFileExists{parskip.sty}{%
    \usepackage{parskip}
  }{% else
    \setlength{\parindent}{0pt}
    \setlength{\parskip}{6pt plus 2pt minus 1pt}}
}{% if KOMA class
  \KOMAoptions{parskip=half}}
\makeatother
\usepackage{xcolor}
\usepackage[margin=18mm]{geometry}
\usepackage{longtable,booktabs,array}
\usepackage{calc} % for calculating minipage widths
% Correct order of tables after \paragraph or \subparagraph
\usepackage{etoolbox}
\makeatletter
\patchcmd\longtable{\par}{\if@noskipsec\mbox{}\fi\par}{}{}
\makeatother
% Allow footnotes in longtable head/foot
\IfFileExists{footnotehyper.sty}{\usepackage{footnotehyper}}{\usepackage{footnote}}
\makesavenoteenv{longtable}
\setlength{\emergencystretch}{3em} % prevent overfull lines
\providecommand{\tightlist}{%
  \setlength{\itemsep}{0pt}\setlength{\parskip}{0pt}}
\setcounter{secnumdepth}{-\maxdimen} % remove section numbering
% Shared preamble for the pandoc-generated roadmap PDFs (pdflatex).
\usepackage[T1]{fontenc}
\usepackage[utf8]{inputenc}
\usepackage{lmodern}
\usepackage{microtype}
\usepackage{amssymb}
\usepackage{booktabs}
\usepackage{longtable}
\usepackage{array}
\usepackage{xcolor}

% Map the non-ASCII glyphs the markdown uses so pdflatex doesn't choke.
\usepackage{newunicodechar}
\newunicodechar{→}{$\rightarrow$}
\newunicodechar{←}{$\leftarrow$}
\newunicodechar{↔}{$\leftrightarrow$}
\newunicodechar{⇒}{$\Rightarrow$}
\newunicodechar{±}{$\pm$}
\newunicodechar{≈}{$\approx$}
\newunicodechar{≥}{$\geq$}
\newunicodechar{≤}{$\leq$}
\newunicodechar{×}{$\times$}
\newunicodechar{÷}{$\div$}
\newunicodechar{·}{$\cdot$}
\newunicodechar{…}{\dots}
\newunicodechar{²}{\textsuperscript{2}}
\newunicodechar{³}{\textsuperscript{3}}
\newunicodechar{°}{\textdegree}
\newunicodechar{µ}{$\mu$}
\newunicodechar{σ}{$\sigma$}
\newunicodechar{ε}{$\varepsilon$}
\newunicodechar{θ}{$\theta$}
\newunicodechar{τ}{$\tau$}
\newunicodechar{γ}{$\gamma$}
\newunicodechar{ω}{$\omega$}
\newunicodechar{ζ}{$\zeta$}
\newunicodechar{Σ}{$\Sigma$}
\newunicodechar{√}{$\surd$}
\newunicodechar{✓}{\checkmark}
\newunicodechar{•}{\textbullet}
\newunicodechar{≠}{$\neq$}
\newunicodechar{∞}{$\infty$}
\newunicodechar{ }{~}
% status emoji + checks -> short text (pdflatex has no emoji glyphs)
\newunicodechar{🔬}{{[study]}}
\newunicodechar{🔴}{{[planned]}}
\newunicodechar{🟡}{{[planned]}}
\newunicodechar{🟢}{{[easy]}}
\newunicodechar{🗺}{{[map]}}
\newunicodechar{✅}{{[done]}}
\newunicodechar{⚠}{{[!]}}
\newunicodechar{️}{}   % U+FE0F variation selector
\newunicodechar{🤖}{}
\ifLuaTeX
  \usepackage{selnolig}  % disable illegal ligatures
\fi
\IfFileExists{bookmark.sty}{\usepackage{bookmark}}{\usepackage{hyperref}}
\IfFileExists{xurl.sty}{\usepackage{xurl}}{} % add URL line breaks if available
\urlstyle{same}
\hypersetup{
  colorlinks=true,
  linkcolor={blue},
  filecolor={Maroon},
  citecolor={Blue},
  urlcolor={blue},
  pdfcreator={LaTeX via pandoc}}

\author{}
\date{}

\begin{document}

\hypertarget{prior-art-mechanism-reference-automated-multirotor-control-tuning}{%
\section{Prior-art mechanism reference: automated multirotor control
tuning}\label{prior-art-mechanism-reference-automated-multirotor-control-tuning}}

Companion to \url{gcs-live-rig-tuning.md}. A builder\textquotesingle s
teardown of every comparable system found in the
\href{gcs-live-rig-tuning.md\#feasibility-verdict}{deep-research
survey}: how each one is \textbf{set up}, the \textbf{algorithm}, the
\textbf{safety system}, how it checks \textbf{correctness}, its
\textbf{communication}, and the \textbf{control mechanics} (where and
how it excites the plant). Written so an engineer can lift the right
pieces into their own tuner.

\begin{quote}
This file is the \textbf{at-a-glance} view. For the \textbf{from-core}
treatment of each system --- control-theory background, block diagrams,
algorithm flowcharts, protocol sequences and signal sketches --- see
\url{gcs-live-rig-tuning-prior-art-deep.md}.
\end{quote}

\begin{quote}
\textbf{Confidence / currency.} Architecture-level facts (where the
optimizer runs, excitation type, in-flight vs offline, openness) are
from primary docs/papers and were adversarially verified (see the
study). Lower-level algorithm details (model order, RLS, exact cost
terms) are established but \textbf{version-specific} --- PX4/ArduPilot
track their \texttt{main} branch; verify against the current source
before copying constants. Items marked \emph{(general)} are well-known
mechanism, not separately source-verified here.
\end{quote}

\hypertarget{at-a-glance}{%
\subsection{At a glance}\label{at-a-glance}}

\begin{longtable}[]{@{}lllll@{}}
\toprule\noalign{}
System & Optimizer runs & Excitation & Where & Open? \\
\midrule\noalign{}
\endhead
\bottomrule\noalign{}
\endlastfoot
ArduCopter AUTOTUNE & \textbf{onboard FC} (iterative twitch search) &
rate step "twitch", axis-by-axis & \textbf{in free flight} & GPLv3 \\
ArduCopter SystemID & \textbf{offline} (model fit + PID opt) & frequency
chirp on a mixer axis & in free flight (log), fit offline & GPLv3 \\
PX4 \texttt{mc\_autotune} & \textbf{onboard FC} (sys-ID + gain calc) &
per-axis rate/torque disturbance (doublet) & in free flight
(alt-stabilized) & BSD-3 \\
PX4 manual/step guide & \textbf{human} & stick step input & in free
flight & BSD-3 (docs) \\
CIFER + CONDUIT/Simulink (academic) & \textbf{offline desktop} &
frequency sweeps (pilot/autopilot) & in free flight (log), fit offline &
proprietary tools, published method \\
MathWorks Closed-Loop PID Autotuner & \textbf{in the block} (online
FRF→gains) & sinusoidal perturbation at PID \textbf{output} & sim,
deployable to HW; free flight & proprietary \\
Safe Bayesian Opt (Berkenkamp et al.) & \textbf{offline/host optimizer}
picks candidates, evaluates on HW & controller run per candidate
(closed-loop trial) & in free flight & published method \\
Thrust stands (RCbenchmark/Tyto) & n/a (characterization) & motor/prop
sweep on a stand & \textbf{on a bench} (component, not closed-loop PID)
& proprietary HW + scripting \\
\end{longtable}

\textbf{None} runs a \textbf{live GCS optimizer with the FC as a
per-pass cost-reporting executor}, and \textbf{none} tunes on a
\textbf{mechanically-constrained rig} --- those two are ours (see the
study\textquotesingle s novelty section).

\begin{center}\rule{0.5\linewidth}{0.5pt}\end{center}

\hypertarget{1-arducopter-autotune-onboard-in-flight-twitch-search}{%
\subsection{1. ArduCopter AUTOTUNE (onboard, in-flight, twitch
search)}\label{1-arducopter-autotune-onboard-in-flight-twitch-search}}

Source: \url{https://ardupilot.org/copter/docs/autotune.html}

\begin{itemize}
\tightlist
\item
  \textbf{Setup.} Tune one axis-set per flight; needs calm conditions
  and open space (the craft repeatedly twitches and drifts). Operator
  flies to a safe altitude in AltHold/Loiter, flips an RC aux switch (or
  mode) to AUTOTUNE, and \textbf{holds position while keeping a hand on
  the sticks}. An aggressiveness parameter (\texttt{AUTOTUNE\_AGGR},
  \textasciitilde0.05--0.10) trades responsiveness vs overshoot;
  \texttt{AUTOTUNE\_AXES} selects roll/pitch/yaw.
\item
  \textbf{Control mechanics.} Injects a \textbf{rate "twitch"} (a brief
  commanded angular- rate step) directly into the rate controller of the
  axis under test, then watches the gyro response. \emph{(general)}
  Excitation is generated \textbf{onboard}.
\item
  \textbf{Algorithm.} Per-axis iterative search, not a model fit: ramp
  \textbf{rate-D} up until the response just starts to
  oscillate/bounce-back, back off; ramp \textbf{rate-P} up to a target
  response; set \textbf{angle-P} from the achieved bandwidth; apply a
  safety margin scaled by \texttt{AUTOTUNE\_AGGR}. Converges
  axis-by-axis over many twitches. \emph{(general)}
\item
  \textbf{Safety.} Original gains are kept; \textbf{moving the sticks
  hands control back} to the pilot instantly; disarming or leaving the
  mode \textbf{restores the pre-tune gains} unless the operator
  explicitly saves them (land + low-throttle save, or switch).
  Self-aborts on excessive attitude error. \emph{(general)}
\item
  \textbf{Correctness.} Convergence is per-axis to a target
  rate-response shape; the proof is the \textbf{post-tune test flight}
  (the pilot flies and decides to save). No model/coherence metric.
\item
  \textbf{Communication.} Self-contained on the FC; the GCS (Mission
  Planner/QGC) only sets params (MAVLink \texttt{PARAM\_SET}) and arms
  the mode. Results are the updated PID params, saved to the FC.
\item
  \textbf{Takeaway for us.} The \emph{twitch + bounded per-axis search}
  is the spiritual ancestor of our doublet rollout --- but the search
  lives on the FC. We move that search to the GCS. Their
  \textbf{stick-override abort} maps onto our operator gesture.
\end{itemize}

\hypertarget{2-arducopter-systemid-chirp-excitation--offline-modelpid-fit}{%
\subsection{2. ArduCopter SystemID (chirp excitation + offline model/PID
fit)}\label{2-arducopter-systemid-chirp-excitation--offline-modelpid-fit}}

Sources:
\url{https://ardupilot.org/copter/docs/systemid-model-development.html},
\url{https://ardupilot.org/copter/docs/common-systemid-mode-operation.html}

\begin{itemize}
\tightlist
\item
  \textbf{Setup.} A dedicated \texttt{SystemID} flight mode. The
  operator hovers and triggers a sweep; logs are pulled afterward for
  desktop analysis. Params pick the axis and magnitude.
\item
  \textbf{Control mechanics.} Injects a \textbf{frequency-sweep (chirp)}
  by \emph{superimposing the waveform on the mixer inputs} via
  \texttt{SID\_AXIS} (10 = roll, 11 = pitch, 12 = yaw) --- "direct input
  to the regarded system." The chirp is generated onboard
  (\texttt{chirp\_input.init()/update()},
  \texttt{actuator\_roll\_sysid()} in \texttt{mode\_systemid.cpp}).
\item
  \textbf{Algorithm.} Two clean stages: (1) \textbf{offline}
  frequency-domain fit --- plant transfer-function parameters are
  optimized to best match the collected flight- data frequency
  responses; (2) the identified model is then used to \textbf{optimize
  the PID} in a separate step. Identification is decoupled from gain
  tuning.
\item
  \textbf{Safety.} Short bounded sweeps in a stable hover; pilot retains
  override.
\item
  \textbf{Correctness.} Goodness-of-fit of the transfer function to the
  measured FRF; then verify the optimized PID in flight.
\item
  \textbf{Communication.} Onboard logging (dataflash) → offline tools on
  a PC; gains uploaded back as params.
\item
  \textbf{Takeaway for us.} Confirms \textbf{FC-side chirp at the
  mixer/setpoint level} is a proven excitation injection point ---
  exactly our "override the axis setpoint." The clean
  \textbf{ID-then-tune} split is an alternative to our direct cost
  search; we chose direct cost (no model), but their injection mechanics
  transfer directly.
\end{itemize}

\hypertarget{3-px4-mc_autotune-fully-onboard-sys-id--gain-calc-in-flight}{%
\subsection{\texorpdfstring{3. PX4 \texttt{mc\_autotune} (fully onboard
sys-ID + gain calc,
in-flight)}{3. PX4 mc\_autotune (fully onboard sys-ID + gain calc, in-flight)}}\label{3-px4-mc_autotune-fully-onboard-sys-id--gain-calc-in-flight}}

Sources: \url{https://docs.px4.io/main/en/config/autotune_mc},
\url{https://docs.px4.io/main/en/msg_docs/AutotuneAttitudeControlStatus}

\begin{itemize}
\tightlist
\item
  \textbf{Setup.} Vehicle must be \textbf{flying in an
  altitude-stabilized mode} (not a bench). Operator triggers autotune
  from QGroundControl; it runs hands-off.
\item
  \textbf{Control mechanics.} The flight stack \textbf{applies a small
  disturbance to each axis} (a normalized torque/rate setpoint
  perturbation; effectively a doublet), recording the filtered input
  \texttt{u\_filt} (normalized torque setpoint) and output
  \texttt{y\_filt} (angular velocity).
\item
  \textbf{Algorithm.} \textbf{Onboard system identification}: fit a
  discrete-time SISO model (a low-order 2-pole/2-zero ARX, via recursive
  least squares \emph{(general/version- specific)}) --- the
  \texttt{AutotuneAttitudeControlStatus} uORB message carries
  \texttt{coeff} (model coefficients), \texttt{coeff\_var} (variance),
  and a \textbf{fitness}. From the identified model it computes the
  gains directly: \texttt{kc,\ ki,\ kd,\ kff,\ att\_p}. Processes axes
  \textbf{sequentially roll → pitch → yaw}.
\item
  \textbf{Safety.} \textbf{Reverts unstable gains in real time onboard};
  requires the altitude-stabilized envelope; operator can abort by
  switching mode. The disturbance amplitude is deliberately small.
\item
  \textbf{Correctness.} The onboard \textbf{fitness} + coefficient
  \textbf{variance} gate acceptance; final check is a verification
  flight.
\item
  \textbf{Communication.} Triggered over \textbf{MAVLink} from QGC;
  status streamed via the uORB→MAVLink bridge; resulting params written
  to the FC.
\item
  \textbf{Takeaway for us.} The strongest precedent for \textbf{FC-side
  excitation + onboard scoring} --- PX4 already computes a fitness on
  the FC, which is our "FC computes cost." The difference: PX4
  \emph{also} computes the gains onboard from a model; we keep the
  \textbf{iterative gain choice on the GCS} and skip the model (direct
  cost).
\end{itemize}

\hypertarget{4-px4-manual--step-input-tuning-guide-human-in-the-loop-baseline}{%
\subsection{4. PX4 manual / step-input tuning guide (human-in-the-loop
baseline)}\label{4-px4-manual--step-input-tuning-guide-human-in-the-loop-baseline}}

Source:
\url{https://docs.px4.io/main/en/config_mc/pid_tuning_guide_multicopter}

\begin{itemize}
\tightlist
\item
  \textbf{Setup / mechanics / algorithm.} No optimizer. Hover,
  \textbf{give a fast step input via the stick} ("push roll to one side,
  then let it snap back"), watch for oscillation/overshoot, \textbf{land
  before changing a parameter}, repeat.
\item
  \textbf{Safety.} The land-between-changes rule is the safety system;
  slow throttle ramps to check for oscillation first.
\item
  \textbf{Takeaway.} This is the manual loop our autotuner automates;
  the \textbf{step doublet} is the same excitation, and "land between
  changes" is what our \textbf{operator-gated,
  re-stabilize-between-rollouts} protocol replaces with a rig.
\end{itemize}

\hypertarget{5-cifer--conduit--simulink-academic-offline-frequency-domain}{%
\subsection{5. CIFER + CONDUIT / Simulink (academic offline
frequency-domain)}\label{5-cifer--conduit--simulink-academic-offline-frequency-domain}}

Sources:
\url{https://www.academia.edu/54509862/Frequency_Response_System_Identification_and_Flight_Controller_Tuning_for_Quadcopter_UAV},
\url{https://arxiv.org/pdf/1508.04886},
\url{https://www.sjsu.edu/researchfoundation/docs/AHS_2015_Wei.pdf},
\url{https://arxiv.org/pdf/0804.3881}

\begin{itemize}
\tightlist
\item
  \textbf{Setup.} Instrument the craft, fly \textbf{frequency sweeps in
  hover}, axis by axis. Sweeps are input by a \textbf{human pilot via
  RC}, or (better) \textbf{generated by the autopilot} --- a human
  "cannot accurately cover the whole frequency needed," so an
  onboard-generated sweep yields higher-quality data (Adiprawita et
  al.).
\item
  \textbf{Control mechanics.} Frequency-sweep excitation per axis,
  logged.
\item
  \textbf{Algorithm.} \textbf{CIFER} (Comprehensive Identification from
  FREquency Response) does multi-window FFT → frequency response → fits
  SISO transfer functions in the \textbf{frequency domain} (offline).
  Then PID gains are \textbf{optimized offline} against
  handling-qualities specs with \textbf{CONDUIT} or \textbf{MATLAB
  Simulink Design Optimization}. Gains are uploaded to the Pixhawk
  (overwritten via QGC).
\item
  \textbf{Safety.} Conventional piloted flight with a safety pilot; the
  heavy lifting is offline, so nothing risky happens autonomously on the
  vehicle.
\item
  \textbf{Correctness.} The \textbf{coherence function} validates each
  frequency point of the FRF (data quality), and the model fit cost +
  the CONDUIT spec margins gauge the result; verified by flight test.
\item
  \textbf{Communication.} Onboard log → desktop tools → params back via
  QGC/MAVLink. Pure \textbf{upload-conduit} GCS role.
\item
  \textbf{Takeaway for us.} The rigorous end of the field.
  \textbf{Coherence} is a correctness idea worth stealing --- a
  per-window data-quality gate before trusting a cost. But
  it\textquotesingle s offline and model-based; we trade model fidelity
  for closed-loop directness.
\end{itemize}

\hypertarget{6-mathworks-closed-loop-pid-autotuner}{%
\subsection{6. MathWorks Closed-Loop PID
Autotuner}\label{6-mathworks-closed-loop-pid-autotuner}}

Sources:
\url{https://www.mathworks.com/help/slcontrol/ug/pid-controller-tuning-for-a-uav-quadcopter.html},
\url{https://www.mathworks.com/help/slcontrol/ug/closedlooppidautotuner.html},
\url{https://www.mathworks.com/help/uav/ug/pid-autotuning-for-uav-quadcopter.html}

\begin{itemize}
\tightlist
\item
  \textbf{Setup.} A Simulink block in the control loop; primarily
  demonstrated in simulation, but \textbf{deployable to real hardware}
  via generated C/C++ ("use the Closed-Loop PID Autotuner on hardware to
  perform the same process").
\item
  \textbf{Control mechanics.} Injects \textbf{sinusoidal perturbation
  signals at the OUTPUT of each existing PID controller} (8 loops in the
  quad example) --- i.e. it perturbs the actuator command, not the
  setpoint.
\item
  \textbf{Algorithm.} Estimates the plant \textbf{frequency response in
  closed loop, in real time}, from the perturbation experiment, then
  computes PID gains from the FRF. Self-contained (the optimizer is in
  the block, not external).
\item
  \textbf{Safety.} Runs the experiment around the operating point with
  bounded perturbations; closed-loop so the existing controller keeps
  the craft stable during the experiment.
\item
  \textbf{Correctness.} FRF-based; the block reports estimated margins.
\item
  \textbf{Takeaway for us.} The "\textbf{perturb at the controller
  output, identify in closed loop}" trick is an alternative excitation
  injection point (actuator vs setpoint) and avoids needing clean
  setpoint authority --- worth noting if setpoint override proves
  awkward on the FC.
\end{itemize}

\hypertarget{7-safe-bayesian-optimization-on-hardware-berkenkamp-et-al}{%
\subsection{7. Safe Bayesian Optimization on hardware (Berkenkamp et
al.)}\label{7-safe-bayesian-optimization-on-hardware-berkenkamp-et-al}}

Source: \url{https://las.inf.ethz.ch/files/berkenkamp16safe.pdf}

\begin{itemize}
\tightlist
\item
  \textbf{Why it matters most to us.} This is the \textbf{closest
  architectural cousin}: an \textbf{external optimizer proposes
  candidate controller parameters, each is evaluated by a real
  closed-loop trial on the quadrotor, a scalar performance is measured,
  and the optimizer picks the next candidate} --- exactly our
  GCS-as-optimizer / evaluate-a-pass loop, minus the rig.
\item
  \textbf{Algorithm.} \textbf{Safe Bayesian Optimization (SafeOpt)}: a
  Gaussian-process model of performance-vs-parameters with a
  \textbf{safety constraint}, so the search only tries parameter sets
  predicted to stay above a safety threshold --- it never knowingly
  evaluates a destabilizing gain. Sample-efficient (few trials).
\item
  \textbf{Safety / correctness.} The safety constraint is the headline
  feature: it bounds exploration to a "safe set," directly addressing
  our "the search probes unstable gains by design" risk.
\item
  \textbf{Takeaway for us.} If our existing SITL optimizer struggles
  with real noise (the study\textquotesingle s deferred "optimizer
  selection" question), \textbf{SafeOpt-style Bayesian optimization is
  the strongest candidate} --- sample-efficient (few real rollouts) and
  safety-constrained (won\textquotesingle t probe known-bad gains).
  Strongly consider for the live phase even though we keep the rig as
  the hard safety net.
\end{itemize}

\hypertarget{8-adjacent-not-closed-loop-pid-tuners-but-useful-inputs}{%
\subsection{8. Adjacent (not closed-loop PID tuners, but useful
inputs)}\label{8-adjacent-not-closed-loop-pid-tuners-but-useful-inputs}}

\begin{itemize}
\tightlist
\item
  \textbf{Thrust stands --- RCbenchmark / Tyto Robotics}
  (\url{https://www.tytorobotics.com/pages/rcbenchmark-software},
  \url{https://www.tytorobotics.com/blogs/software-troubleshooting/automatic-control-pid}):
  bench rigs that \textbf{characterize a single motor+prop} (thrust,
  torque, current, efficiency, and a PID for the
  \emph{stand\textquotesingle s own} control). \textbf{Not} closed-loop
  airframe PID tuning. Use them to \textbf{feed real motor/prop
  constants into the SITL model}, improving sim-then-hardware transfer
  --- not to tune the FC.
\item
  \textbf{Data-driven parameter ID + RL} (Eschmann/Albani/Loianno 2024,
  \url{https://arxiv.org/pdf/2404.07837}): identifies
  inertia/thrust/motor-delay from \textasciitilde73 s of \textbf{free
  flight}, then trains an RL policy. Confirms free-flight data
  collection dominates; a pointer for a future model-based or learned
  controller, not classical PID.
\item
  \textbf{Betaflight}: no built-in closed-loop PID autotuner
  \emph{(general, not source-verified here)} --- relies on presets,
  filters, and manual tuning. The research found no confirmed Betaflight
  autotune claim.
\end{itemize}

\begin{center}\rule{0.5\linewidth}{0.5pt}\end{center}

\hypertarget{synthesis--what-to-borrow-mapped-to-our-subsystems}{%
\subsection{Synthesis --- what to borrow, mapped to our
subsystems}\label{synthesis--what-to-borrow-mapped-to-our-subsystems}}

\begin{longtable}[]{@{}ll@{}}
\toprule\noalign{}
Our subsystem & Best precedent to copy \\
\midrule\noalign{}
\endhead
\bottomrule\noalign{}
\endlastfoot
FC-side excitation injection point & ArduCopter SystemID (chirp on the
mixer/setpoint axis); MathWorks (perturb at PID output) as a fallback \\
Excitation waveform & PX4 doublet + ArduPilot/CIFER chirp --- support
both (we already do) \\
Onboard cost/score & PX4 onboard \textbf{fitness} + coefficient
variance \\
Data-quality gate & CIFER \textbf{coherence} per window → reject bad
rollouts before scoring \\
Optimizer (if SITL one underperforms) & \textbf{SafeOpt / Bayesian
optimization} (sample-efficient + safety-constrained) \\
Abort / handback & ArduPilot AUTOTUNE \textbf{stick-override} → our
operator gesture \\
Keep-old-gains-until-saved & ArduPilot AUTOTUNE / PX4 revert → our
"never persist intermediates; explicit Apply" \\
Sim-then-hardware & Standard across PX4/ArduPilot; feed real motor
constants from a thrust stand \\
\end{longtable}

The two things \textbf{no precedent gives us} --- a \textbf{live GCS
optimizer with the FC as a per-pass executor}, and a
\textbf{constrained-rig closed-loop rollout} --- are where the design
work is genuinely original, and where the study\textquotesingle s open
risks (rig-induced dynamics bias, reliable command/result over NavLink
v2) have no off-the-shelf answer to copy.

\hypertarget{changelog}{%
\subsection{Changelog}\label{changelog}}

\begin{longtable}[]{@{}lll@{}}
\toprule\noalign{}
Date & Author & Description \\
\midrule\noalign{}
\endhead
\bottomrule\noalign{}
\endlastfoot
15/06/2026 & ragnar-vallhala & Initial prior-art mechanism reference \\
\end{longtable}

\end{document}
